\let\cleardoublepage\clearpage
\addcontentsline{toc}{chapter}{Abstract}
\renewcommand{\abstractname}{Abstract}


\thispagestyle{plain}
\begin{center}
    \textbf{ECONOMIC TIME SERIES MODELING: FROM THE RELATIVE CHANGE RATE TO STOCHASTICALLY STRUCTURED REGIME TRANSITION MODELS}

        \vspace{0.9cm}
    \textbf{Abstract}
\end{center}

Regime analysis in economic and financial time series presents a persistent trade-off between deterministic models, which lack stochastic flexibility, and probabilistic models that are difficult to interpret.  This research proposes a unified framework to bridge that gap, articulated in three progressive contributions.

The first contribution is the formalization of the Optimal Commuted Relative Change Rate (TCROC) operator family and its parametric variants (TCROCM, ETCROC, ETCROCM), with rigorous proofs of existence, uniqueness, asymptotic consistency, and perturbation stability.  These operators transform a continuous time series into a sequence of interpretable economic regimes.

The second contribution integrates the TCROC operator with Markov chains and estimates the transition matrix using a Non-Negative Least Squares (NNLS) algorithm.  Unlike the maximum likelihood approach used in Markov-Switching (MS-AR) models, the NNLS formulation is convex, guarantees global convergence, and ensures stochastic validity of the estimated matrix by construction.

The third contribution extends the linear framework through Stochastically Structured Reservoir Computing (SSRC).  It is formally proven that the TCROC-Markov model is a special case of SSRC (Inclusion Theorem), providing a theoretical path toward capturing nonlinear dynamics.  Empirical validation of this extension shows 20--23\% improvements in forecast accuracy over the linear model.

Comprehensive empirical validation is carried out on weekly series of four Honduran fuels (Regular, Super, Diesel, and Kerosene) during 2017--2025, including exhaustive hyperparameter optimization, spectral analysis, and comparison with benchmarks (fixed quantiles and MS-AR).  The results confirm the statistically significant superiority of the TCROC-Markov model and reveal its dual utility: high predictive accuracy with extensive data and dynamic structure capture with limited data.

\textbf{Keywords:} TCROC, time series, Markov chains, NNLS, regime switching, reservoir computing, SSRC, fuel prices, Honduras.
